\documentclass[10pt,a4paper]{amsart}
\usepackage[margin=19mm]{geometry}
\usepackage{amsmath,amssymb,mathtools}
\usepackage[scr=rsfs,cal=euler]{mathalfa}
\usepackage{xfrac}
\usepackage[unicode,hidelinks,bookmarks=false]{hyperref}
\newcommand{\ie}{i.e.\ }
\newcommand{\cf}{cf.\ }
\newcommand{\resp}{respectively\ }
\newcommand{\Subsection}{\subsection}
\providecommand{\nobreakdash}{\nobreak\hbox{-}}
\setlength{\emergencystretch}{2em}
\multlinegap0pt
\usepackage{bm}
\usepackage{bbm}
\usepackage{dsfont}
\usepackage{xspace}
\usepackage{cancel}
\usepackage{mathtools}
\usepackage{amsthm}
\makeatletter
\@ifundefined{lemma}{\newtheorem{lemma}{Lemma}}{}
\@ifundefined{proposition}{\newtheorem{proposition}{Proposition}}{}
\@ifundefined{theorem}{\newtheorem{theorem}{Theorem}}{}
\makeatother
\newcommand{\Z}{\mathbb Z}
\newcommand{\id}{\mathrm{id}}
\newcommand{\one}{\mathbbm{1}}
\title[No quantum solutions to linear constraint systems from monom]{No quantum solutions to linear constraint systems from monomial measurement-based quantum computation in odd prime dimension}
\author{Markus Frembs, Cihan Okay, Ho Yiu Chung}
\date{Source: arXiv:2601.11367v1 (CC BY 4.0)}
\begin{document}
\maketitle

\noindent\emph{Source and licence.} No quantum solutions to linear constraint systems from monomial measurement-based quantum computation in odd prime dimension. Version arXiv:2601.11367v1 (CC BY 4.0), distributed under the Creative Commons Attribution 4.0 International licence, as recorded in the arXiv licence metadata for this exact version and shown on the abstract page https://arxiv.org/abs/2601.11367v1. Full rights notice: licenses/arxiv-2601.11367-CC-BY-4.0.txt.\par\smallskip

\noindent\textit{Editorial scope.} Counted unit: the Theorem whose source label is \texttt{thm: main L to H} in the article (source0.tex lines 989--992, source label \texttt{thm: main L to H}), delivered together with the article's own complete original proof of it (source0.tex lines 994--1034, one \texttt{proof} environment of 41 lines). No mathematics is written, repaired or re-derived: statement and proof are the article's own text line for line. Required context retained uncounted: lem: comm\_tildeK (lines 760--768); thm:K to H (lines 918--920); lem: action invariant (lines 941--948); prop:tp to h, homo, z\_hat invariant (lines 962--964). Every definition, display and label of those lines is retained unchanged. Omitted from this package and not redistributed: 1--759, 769--917, 921--940, 949--961, 965--988, 993--993, 1035--1226. Nothing inside a retained excerpt is omitted or altered: every retained body is delivered exactly as the article's lines are written, so no omission receipt of any kind accompanies this package. Citations: the retained excerpts carry no citation command at all, so no bibliography entry of the article is transcribed and no reference is redistributed. Reference adaptation: none was needed and none was made; every reference inside the delivered unit resolves inside it, so no unresolved reference or citation is produced. Printed slips kept literal: no printed word, symbol, hypothesis or number of the counted statement or of its proof was altered. Layout and class: the article's own class is \texttt{revtex4-1} with packages including amsfonts, amsmath, amssymb, amsthm, mathtools, blkarray, nicematrix, enumerate, bm, dsfont, graphicx, nicefrac, color, xy; the wrapper is the lane's pinned \texttt{amsart} editorial wrapper at 10pt on A4, which restates the theorem-like environments the retained text needs and adds the article's own macro definitions that the retained text needs (\texttt{\textbackslash Z}, \texttt{\textbackslash id}, \texttt{\textbackslash one}), with extra wrapper packages (bm, bbm, dsfont, xspace, cancel, mathtools). The delivered numbering is the unit's own. Assets: no retained span contains a figure environment, a graphics-inclusion command, a PDF-page inclusion, a table or a TikZ picture; no image file is copied into this package, the package declares no asset dependency and no image file is redistributed with it. Licence: version arXiv:2601.11367v1 of this article is distributed under the Creative Commons Attribution 4.0 International licence, as recorded in the arXiv licence metadata for this exact version and shown on the abstract page https://arxiv.org/abs/2601.11367v1. A full rights notice is delivered at licenses/arxiv-2601.11367-CC-BY-4.0.txt. Characters: every retained excerpt is pure ASCII, so the delivered engine, XeLaTeX with its default Latin Modern text font, reports no missing character.
\begin{lemma}\label{lem: comm_tildeK}
    Let $M=(\xi,\sigma),M'=(\xi',\sigma')\in N$ with $M^p,M'^p=\one$ and $[M,M']\in \langle\omega\one\rangle=Z(N)$ for $\omega=e^{\frac{2\pi i}{p}}$.
    \begin{enumerate}
    \item If $\sigma=\sigma'=\id$, then $M,M'\in T$.
    \item If $\sigma=\id$ and $\sigma'\neq\id$, then $\phi^N_{S_p}(M')\cdot M\in H\cap T_{(p)}$. In particular, if $[M,M']=\one$, then $M\in Z(N)$.
    \item If $\sigma,\sigma'\neq\id$ then $\phi^N_{S_p}(M)=\phi^N_{S_p}(M')$ and there exists $a\in\Z_p$ and $(\chi,\id)\in H\cap T_{(p)}$ such that
    $$\left(\phi^N_{S_p}(M')\cdot M'\right)=\left((\phi^N_{S_p}(M)\cdot M)(\chi,\id)\right)^a$$ In particular, if $[M,M']=\one$, then $(\chi,\id)\in Z(N)$.
    \end{enumerate}
\end{lemma}

\begin{theorem}{\label{thm:K to H}}
    Let $M,M'\in K$ such that $M^p=M'^p=\one$ and $[M,M']\in Z(K)$, then $\phi^K_H(MM')=\phi^K_H(M)\phi^K_H(M')$.
\end{theorem}

\begin{lemma}{\label{lem: action invariant}}
    Let $M\in T_{(p)}$ and $\tau\in Im(\phi^N_{S_p})$. Then the following relations hold:
    \begin{enumerate}
        \item $\phi^{T_{(p)}}_H(\tau\cdot M)=\phi^{T_{(p)}}_H(M)$
        \item $\phi^L_H(M)=\phi^{T_{(p)}}_H(M)=\phi^K_H(M)$
        \item $\phi^L_H(\tau\cdot M)=\phi^L_H( M)$
    \end{enumerate}
\end{lemma}

\begin{proposition}{\label{prop:tp to h, homo, z_hat invariant}}
    $\phi^{T_{(p)}}_H$ is a homomorphism and $\phi^{T_{(p)}}_H(M)=M$ for all $M\in H\cap T_{(p)}$.
\end{proposition}

\begin{theorem}{\label{thm: main L to H}}
    Let $M,M'\in L$ such that $[M,M']\in Z(K)$ and $M^p_1=M^p_2=\one$, then
    $$\phi^L_H(MM')=\phi^L_H(M)\phi^L_H(M')\; .$$
\end{theorem}

\begin{proof}
    We have the following cases to consider:
    \begin{enumerate}
    \item $M,M'\in T_{(p)}$
    \item $M\in T_{(p)}$ and $M'\not\in T_{(p)}$ (analogously, $M\not\in T_{(p)}$ and $M'\in T_{(p)}$)
    \item $M,M'\not\in T_{(p)}$
    \end{enumerate}
    (1): Since $M,M'\in T_{(p)}$, also $MM'\in T_{(p)}$. Hence, by Lm.~\ref{lem: action invariant} and Prop.~\ref{prop:tp to h, homo, z_hat invariant}, we find
    \begin{align*}
        \phi^L_H(MM')
        =\phi^{T_{(p)}}_H(MM')
        =\phi^{T_{(p)}}_H(M)\phi^{T_{(p)}}_H(M')
        =\phi^L_H(M)\phi^L_H(M')\; .
    \end{align*}
    (2): Let $\tau=\phi^N_{S_p}(M')$ such that $\tau\cdot M'\in H$. By Lm.~\ref{lem: comm_tildeK}, also $\tau\cdot M\in H\cap T_{(p)}\leq H$. By Lm.~\ref{lem: action invariant} and Thm.~\ref{thm:K to H},
    \begin{align*}
        \phi^L_H(MM')
        &=\phi^K_H(\tau\cdot (MM'))\\
        &=\phi^K_H((\tau\cdot M)(\tau\cdot M'))\\
        &=\phi^K_H(\tau\cdot M)\phi^K_H(\tau\cdot M')\\
        &=\phi^{T_{(p)}}_H(\tau\cdot M)\phi^L_H(M')\\
        &=\phi^{T_{(p)}}_H(M)\phi^L_H(M')\\
        &=\phi^L_H(M)\phi^L_H(M')\; .
    \end{align*}
    (3): By Lm.~\ref{lem: comm_tildeK}, $\phi^N_{S_p}(M)=\phi^N_{S_p}(M')=:\tau$. If $MM'\not\in T_{(p)}$, then also $\tau=\phi^N_{S_p}(MM')$, hence,
    \begin{align*}
        \phi^L_H(MM')
        =\phi^K_H(\tau\cdot(MM'))=\phi^K_H((\tau\cdot M)(\tau\cdot M'))
        =\phi^K_H(\tau\cdot M)\phi^K_H(\tau\cdot M')
        =\phi^L_H(M)\phi^L_H(M')\; .
    \end{align*}
    If $MM'\in T_{(p)}$, then applying Lm.~\ref{lem: action invariant} (1) and Thm.~\ref{thm:K to H} again, we obtain
    \begin{align*}
        \phi^L_H(MN)
        =\phi^L_H(\tau\cdot (MN))
        =\phi^K_H(\tau\cdot (MN))
        &=\phi^K_H((\tau\cdot M)(\tau\cdot N))\\
        &=\phi^K_H(\tau\cdot M)\phi^K_H(\tau\cdot N)
        =\phi^L_H(M)\phi^L_H(N)\; .\qedhere
    \end{align*}
\end{proof}

\end{document}
